%==============================================================================%
%                PRØVE | FORKURS 1P-2P LÆRERUTDANNING | V2017                  %
%==============================================================================%
%
% __/\\\\\\\\\\\\____________________/\\\\\\______________/\\\\\\\\\_____
%  _\/\\\////////\\\_________________\////\\\____________/\\\///////\\\___
%   _\/\\\______\//\\\___________________\/\\\___________\///______\//\\\__
%    _\/\\\_______\/\\\_____/\\\\\\\\_____\/\\\_____________________/\\\/___
%     _\/\\\_______\/\\\___/\\\/////\\\____\/\\\__________________/\\\//_____
%      _\/\\\_______\/\\\__/\\\\\\\\\\\_____\/\\\_______________/\\\//________
%       _\/\\\_______/\\\__\//\\///////______\/\\\_____________/\\\/___________
%        _\/\\\\\\\\\\\\/____\//\\\\\\\\\\__/\\\\\\\\\_________/\\\\\\\\\\\\\\\_
%         _\////////////_______\//////////__\/////////_________\///////////////_
%
%==============================================================================%
%                               MED HJELPEMIDLER                               %
%==============================================================================%

\Del*{m}


%==============================================================================%
%                                 OPPGAVE 2.1                                  %
%==============================================================================%
\Oppgave[2] \points*{2}

Siri kjøper $\SI{1,5}{\kg}$ epler i butikk $A$. Full pris for eplene er $18$
kroner per kilogram.  Siri får $\SI{15}{\percent}$ avslag på denne prisen.
\medskip

Eivind kjøper $\SI{1,5}{\kg}$ epler i butikk $B$. Han får $\SI{10}{\percent}$
avslag på prisen. \medskip

Hva må full pris per kilogram epler være i butikk $B$ for at Eivind skal betale
det samme i butikk $B$ som Siri gjør i butikk $A$?


%==============================================================================%
%                                 OPPGAVE 2.2                                  %
%==============================================================================%
\Oppgave[4] % Oppgave 2.2

\begin{figure}[H]
  \centering
  \adjustbox{valign=t}{
    \begin{subfigure}[t]{0.45\textwidth}
      \centering
      \resizebox{0.4\textwidth}{!}{%
        \tikzsetnextfilename{Forkurs-1p-2p-laererutdanning-2017-V-oppgave-2-2a}
        \begin{tikzpicture}
          \definecolor{brownMedium}{HTML}{623d22}
          \definecolor{brownLight}{HTML}{a78d7b}
          \definecolor{brownDark}{HTML}{573722}
          %stick
          \fill[brownLight]
          (-1,-3.17) -- (-1,-8.48) arc (-140:-39:1.2) -- (0.8,-3.16) --cycle;
          %shadow
          \fill[brownLight!95!black,rounded corners]
          (0.8,-3.15) -- (0.8,-3.9) -- (0.5,-4.3) -- (-0.8,-4.3) -- (-1,-4.1) -- (-1,-3.1);
          %ice cream
          \fill[brownDark]
          (-4.2,5.2) -- (-4.2,-3.1) -- (-3.96,-3.24) -- (3.26,-3.21) -- (3.86,-2.9) -- (3.84,5.8) .. controls (3.54,9.77) and (-3.4,10.05) .. cycle;
          \fill[brownMedium]
          (3.21,5.5) -- (3.26,-3.21) -- (-3.96,-3.24) -- (-4,5.5) .. controls (-2.2,10.1) and (2.41,9.1) .. cycle;
      \end{tikzpicture}}
  \end{subfigure}}%
  ~
  \adjustbox{valign=t}{%
    \begin{subfigure}[t]{0.45\textwidth}
      \centering
      \vspace{0.4cm}
      \tikzsetnextfilename{Forkurs-1p-2p-laererutdanning-2017-V-oppgave-2-2b}
      \begin{tikzpicture}
        \edef\w{2.25}
        \draw[draw=maincolorMedium, thick] (0,\w) -- (0,0) -- (\w,0) node [midway, below]{$\SI{6}{\cm}$} -- (\w,\w) node [midway, right]{$\SI{6}{\cm}$};
        \draw[draw=maincolorMedium, dashed, thick] (0,\w) -- (\w,\w);
        \draw[thick, maincolorMedium] (\w,\w) arc(0:180:0.5*\w);
      \end{tikzpicture}
  \end{subfigure}}
  \caption{}
  \label{fig:Forkurs-1p-2p-laererutdanning-2017-V-oppgave-2-2}
\end{figure}

En pinneis er $\SI{1.5}{cm}$ tykk. Forsiden er et kvadrat med sider
$\SI{6.0}{cm}$ og en halvsirkel. Se
\cref{fig:Forkurs-1p-2p-laererutdanning-2017-V-oppgave-2-2}. I denne
oppgaven ser vi bort fra pinnen.

\begin{oppgaver}
  \Item{2} Bestem volumet av isen.
\end{oppgaver}

\begin{oppgaver}
  \Item{2} Bestem overflaten av isen.
\end{oppgaver}


%==============================================================================%
%                                 OPPGAVE 2.3                                  %
%==============================================================================%
\Oppgave[4]

Du har $200$ mynter i en eske.

\begin{itemize}
  \item $\cfrac{1}{5}$ av myntene er laget før $1940$. Av disse er
    kobbermynter.
  \item av myntene er laget etter $1940$. Halvparten av disse er kobbermynter.
\end{itemize}

Du tar en mynt tilfeldig fra esken.

\begin{oppgaver}
  \Item{2} Bestem sannsynligheten for at du tar en kobbermynt.
\end{oppgaver}

Anta at du tok en kobbermynt.

\begin{oppgaver}
  \Item{2} Bestem sannsynligheten for at denne mynten er laget før $1940$.
\end{oppgaver}


%==============================================================================%
%                                 OPPGAVE 2.4                                  %
%==============================================================================%
\Oppgave[4] % Oppgave 2.4

\begin{figure}[H]
  \tikzsetnextfilename{Forkurs-1p-2p-laererutdanning-2017-V-oppgave-2-4}
  \begin{tikzpicture}
    % Base values ymin=-0.45,ymax=2.85,xmin=-0.4,xmax=5.5
    \tkzInit[ymin=-0.675,ymax=4.275,xmin=-0.6,xmax=8.25]
    \tkzClip
    \edef\scalePic{1.5}
    \pgfmathsetmacro\AC{\scalePic*4}
    \pgfmathsetmacro\BC{\scalePic*3}
    \pgfmathsetmacro\AB{sqrt(\AC*\AC+\BC*\BC)}

    \tkzDefPoint(0,0){A}
    \tkzDefPoint(\AB,0){B}

    \tkzInterCC[R](A,\AC cm)(B,\BC cm)
    \tkzGetPoints{C}{C1}

    \tkzDefMidPoint(A,B) \tkzGetPoint{D}

    \tkzDefPointWith[orthogonal](D,A) \tkzGetPoint{D1}

    \tkzInterLC(A,C)(D,A) \tkzGetPoints{l1}{l2}
    \tkzDefMidPoint(l1,l2) \tkzGetPoint{E}

    \tkzMarkRightAngle(B,C,A)
    \tkzMarkRightAngle(D,E,A)

    \tkzDrawSegments[dashed,thick,color=maincolorLight](C,D D,E)

    \tkzDrawPolygon[color=maincolorDark](A,B,C)
    \tkzLabelPoint[below left](A){$A$}
    \tkzLabelPoint[below right](B){$B$}
    \tkzLabelPoint[above](C){$C$}
    \tkzLabelPoint[below](D){$D$}
    \tkzLabelPoint[above left](E){$E$}
  \end{tikzpicture}
  \caption{}
  \label{fig:Forkurs-1p-2p-laererutdanning-2017-V-oppgave-2-4}
\end{figure}

Gitt \cref{fig:Forkurs-1p-2p-laererutdanning-2017-V-oppgave-2-4}. Punktet $D$
ligger midt på $AB$, $AC=8$ og $BC=6$.

\begin{oppgaver}
  \Item{2} Forklar at $\Delta ABC$ og $\Delta ADE$ er formlike.
\end{oppgaver}

\begin{oppgaver}
  \Item{2} Vis at arealet av $\Delta ACD$ er lik arealet av $\Delta BCD$.
\end{oppgaver}


%==============================================================================%
%                                 OPPGAVE 2.5                                  %
%==============================================================================%
\Oppgave[3] % Oppgave 2.5

Elisabet hoppet lengde med tilløp åtte ganger. Resultatene ser du i
\cref{table:del-2-oppgave-2.5}. \medskip

\begin{table}[H]
   \caption{}
   \label{table:del-2-oppgave-2.5}
   \begin{tabularx}{\textwidth}{ *{4}Y }
    \hline
        \SI{4.98}{\m} & \SI{5.21}{\m} & \SI{5.28}{\m} & \SI{5.07}{\m} \\
        \SI{5.20}{\m} & \SI{4.74}{\m} & \SI{4.85}{\m} & \SI{5.15}{\m} \\
    \hline
    \end{tabularx}
\end{table}

\begin{oppgaver}
  \Item{1} Bestem gjennomsnittet og standardavviket til hopplengdene.
\end{oppgaver}

Tone hoppet også lengde med tilløp åtte ganger. Standardavviket for hopplengdene
til Tone var $\SI{25}{\cm}$.

\begin{oppgaver}
  \Item{2} Hva kan du ut fra dette si om lengdehoppene til Tone sammenliknet
    med lengdehoppene til Elisabet?
\end{oppgaver}


%==============================================================================%
%                                 OPPGAVE 2.6                                  %
%==============================================================================%
\Oppgave[5] \points*{5}

\textquote{Mat på nett} er et firma hvor kunder kan bestille middag som de får
levert på døra. \medskip

Kundene kan velge mellom tre retter:

\begin{itemize}
  \item Dagens fisk koster $110$ kroner.
  \item Dagens kjøtt koster $120$ kroner.
  \item Dagens pasta koster $75$ kroner.
\end{itemize}

Firmaet gir $\SI{10}{\percent}$ rabatt til gode kunder. Levering koster $50$
kroner for avstander som er mindre enn $\SI{8}{\km}$. For lengre avstander er
prisen $100$ kroner. \medskip

Du skal lage ett regneark som firmaet kan bruke for å registrere en bestilling,
legge inn hvor mange prosent rabatt kunden skal få, og beregne hvor mye kunden
skal betale. Regnearket skal se ut som vist nedenfor. I de hvite cellene skal
firmaet registrere opplysninger når de tar imot en bestilling. I de lilla
cellene skal du lage formler.


%==============================================================================%
%                                 OPPGAVE 2.7                                  %
%==============================================================================%
\Oppgave[9]

\Cref{table:Forkurs-1p-2p-laererutdanning-2017-V-oppgave-2-7} viser vekten til
Tonje $x$ måneder etter fødselen.

\begin{table}[H]
    \caption{}
    \label{table:Forkurs-1p-2p-laererutdanning-2017-V-oppgave-2-7}
    \begin{tabularx}{\textwidth}{|l *{7}{Y}|} \hline \Rowcolor
        Måned $x$      &      1     &      2     &      4     &      6     &      8     &      10    &      12    \\ \hline
        Vekt (\si{\g}) & \num{4660} & \num{5300} & \num{6630} & \num{7530} & \num{8210} & \num{8860} & \num{9410} \\ \hline
    \end{tabularx}
\end{table}

\begin{oppgaver}
    \Item{1} Bruk regresjon til å vise at funksjonen $f$ gitt ved
    %
    \begin{equation*}
        f(x) = 4500 \cdot x^{0.29}
    \end{equation*}
    %
    er en god modell for vekten til Tonje det første leveåret.
\end{oppgaver}

\begin{oppgaver}
  \Item{2} Bruk graftegner til å tegne grafen til $f$ for $1 \leq x \leq 12$
\end{oppgaver}

\begin{oppgaver}
  \Item{2} Betem $f(7)$. Hvilken praktisk informasjon gir dette svaret?
\end{oppgaver}

\begin{oppgaver}
  \Item{2} Bestem den gjennomsnittlige vekstfarten til funksjonen $f$ fra $x =
    1$ til $x = 12$ Hvilken praktisk informasjon gir dette svaret?
\end{oppgaver}

\begin{oppgaver}
  \Item{2} Bestem den momentane vekstfarten til funksjonen $f$ når $x=2$
    Hvilken praktisk informasjon gir dette svaret?
\end{oppgaver}


%==============================================================================%
%                                 OPPGAVE 2.8                                  %
%==============================================================================%
\Oppgave[5]

\begin{figure}[H]
  \centering
  \begin{subfigure}[b]{0.06\textwidth}
    \centering
    \resizebox{\textwidth}{!}{\checkers{1}}
    \caption{}
  \end{subfigure}%
  ~
  \begin{subfigure}[b]{0.18\textwidth}
    \centering
    \resizebox{\textwidth}{!}{\checkers{3}}
    \caption{}
  \end{subfigure}%
  ~
  \begin{subfigure}[b]{0.3\textwidth}
    \centering
    \resizebox{\textwidth}{!}{\checkers{5}}
    \caption{}
  \end{subfigure}%
  ~
  \begin{subfigure}[b]{0.42\textwidth}
    \centering
    \resizebox{\textwidth}{!}{\checkers{7}}
    \caption{}
  \end{subfigure}%
  \caption{}
  \label{fig:Forkurs-1p-2p-laererutdanning-2017-V-oppgave-2.8}
\end{figure}

Tenk deg at du skal lage figurer av svarte og hvite kvadrater som vist i
\cref{fig:Forkurs-1p-2p-laererutdanning-2017-V-oppgave-2.8}.

\begin{oppgaver}
  \Item{2} Skriv av og fyll ut
  \cref{table:Forkurs-1p-2p-laererutdanning-2017-V-oppgave-2.8} nedenfor.
\end{oppgaver}

\begin{table}[H]
  \caption{}
  \label{table:Forkurs-1p-2p-laererutdanning-2017-V-oppgave-2.8}
  \begin{tabularx}{\textwidth}{|*{4}{Y|}}
    \hline
    \Rowcolor {Figur\headerstrut} & {Svarte kvadrat} & {Hvite kvadrat} &
    {Kvadrat total} \\ \hline
     1  &  1             &               &                 \\ \hline
     2  &  5             &               &                 \\ \hline
     3  &                &               &                 \\ \hline
     4  &                &               &                 \\ \hline
    $n$ &                &               &                 \\ \hline
  \end{tabularx}
\end{table}

\begin{oppgaver}
    \Item{3} Hvor mange svarte kvadrater trenger du dersom du skal lage en figur
    med totalt $7225$ kvadrater?
\end{oppgaver}


